\section{Worked problems: clamped blocks, mantissa bits and a scale's range}

\subsection{How often does an MX block clamp?}
\textbf{Problem.} Equation~\eqref{eq:mx-exponent} maps a block's largest magnitude
into the element format's top binade, $[2^{e_{\max}}X, 2^{e_{\max}+1}X)$, but the
format's largest value is $t \cdot 2^{e_{\max}}X$ with $t < 2$: 1.5 for E2M1, 1.75
for E4M3, E5M2 and E3M2, 1.875 for E2M3. If the fractional part of
$\log_2 \max_i \lvert v_i\rvert$ is spread evenly over $[0, 1)$ across blocks,
what fraction of blocks clamps its largest value, and by how much at worst?
Compare with the lab's counts on the real projection.

\textbf{Solution.} A block clamps when its maximum, within the top binade, exceeds
$t$ times the binade's start, which happens for a fraction $1 - \log_2 t$ of
evenly spread logarithms: $1 - \log_2 1.5 = 41.5\%$ for E2M1, $1 - \log_2 1.75 =
19.3\%$ for the three formats topping out at 1.75, and $1 - \log_2 1.875 = 9.3\%$
for E2M3 (derived). The lab counted 38.4\%, 18.4\% and 9.2\% of 98,304 blocks
(measured). The worst clamp takes a value just below the next power of two down
to $t$ times the binade's start: it loses $(2 - t)/2$, a quarter of its value in
E2M1 and an eighth in the 1.75 formats. Every clamped block pays that on its
largest, usually most important, element --- which is why NVFP4 fits its scale to
the block instead, and why one could give MX conversion one binade of headroom
and trade the clamp for a coarser grid (the lab's ``break it'' exercise).

\subsection{Mantissa bits predict the table}
\textbf{Problem.} Rounding to a floating-point grid with $m$ mantissa bits has a
step of $2^{-m}$ relative to the binade, so its relative error scales as
$2^{-m}$, as long as the values stay within the format's normal range. Predict
the ratios of the errors of E5M2 to E4M3 and of E2M1 to E2M3, and compare with
the lab's table.

\textbf{Solution.} E5M2 has two mantissa bits and E4M3 three, so E5M2 should be
twice as bad; E2M1 has one and E2M3 three, so MXFP4 should be four times as bad
as MXFP6 (E2M3) (derived). Measured, with one scale per tensor, E5M2's weight and
output errors are 1.98 and 1.99 times E4M3's; in MX blocks, MXFP4's are 4.15 and
4.24 times MXFP6 (E2M3)'s. The prediction ignores the exponent bits entirely, and
it works because a scaled block barely needs them: the lab's projection has
groups of 32 whose largest weight is only 2.6 times their root-mean-square
(\Chref{weight-quantization}), so almost every weight lands within a few
binades of the top. That is also why E3M2 and E5M2 tie: once the block is
scaled, the three extra exponent bits of E5M2 cover values that never occur.
Choosing a format for weights is choosing mantissa bits; exponent bits are for
tensors whose range cannot be scaled away.

\subsection{When does NVFP4's tensor scale bite?}
\textbf{Problem.} In equation~\eqref{eq:nvfp4}, a block with largest magnitude $a_b$
in a tensor with largest magnitude $a_t$ stores the scale
$S_b = 448\,a_b / a_t$ in E4M3, whose smallest normal value is $2^{-6}$ and
smallest subnormal $2^{-9}$. How small can a block's maximum be, relative to the
tensor's, before its scale loses precision, and before it becomes zero? What
does that mean for weights and for activations?

\textbf{Solution.} The scale stays normal while $448\,a_b/a_t \ge 2^{-6}$, that is
while $a_b \ge 3.5 \times 10^{-5}\,a_t$; below that it is subnormal, with fewer
significant bits, and below half the smallest subnormal, $a_b < 2.2 \times
10^{-6}\,a_t$, it rounds to zero and the whole block decodes as zero (derived).
For weights this never happens in practice: a block's maximum is rarely more
than a few hundred times smaller than the tensor's. For an activation with one
outlier channel a thousand times its neighbours it is still three orders of
magnitude away. The limit only bites when a tensor mixes values that differ by
more than $10^{4}$ --- and then a per-tensor FP32 scale was never going to be
enough. MX formats have no such limit, because E8M0 spans $2^{\pm127}$: they
trade NVFP4's exact fit for range, and the lab shows which trade the weights
of a real model prefer.
